% Part: sets-functions-relations
% Chapter: arithmetization
% Section: checking-details
%
\documentclass[../../../include/open-logic-section]{subfiles}


\begin{document}
	
\olfileid{sfr}{arith}{check}
\olsection{مرتبې حلقې او ميدانونه}

په دې ټول څپرکي کښې مو ادعا وکړه چې ځينې تعريفونه «لکه څنګه چې
بايد» هماغسې چلېږي۔ په دې فني ضميمه کښې به روښانه کړو چې مطلب مو
څه دے، او (لنډه خاکه به ورکړو چې څنګه) وښيو تعريفونه «سم» چلېږي۔
\textbf{د ژباړې د نسخې بشپړتيا يادونه (OLARI-006):} سرچينه د حلقې
قوانين ارزوي، خو تر هغو مخکښې دا نۀ ثابتوي چې د صحيح او ناطقو
عددونو پر ټولګيو تعريف شوې عمليې د استازي له انتخابه خپلواکې دي؛
لاندې محاسبې دا اړين ښه‌تعريف‌والے له مخکښې فرضوي۔

په \olref[int]{sec} کښې مو پر $\Int$ جمع او ضرب تعريف کړل۔ غواړو
وښيو چې دا عمليې، لکه څنګه چې تعريف شوې دي، $\Int$ ته هغه جوړښت
ورکوي چې موږ ئې «غواړو»۔ په ځانګړي ډول دا جوړښت يوه تبادلي حلقه ده۔

\begin{defn}
	يوه \emph{تبادلي حلقه} يو سټ $S$ دے چې ځانګړي غړي $0$ او $1$ او عمليې $+$ او $\times$ لري، او دا اته فارمولونه پوره کوي:
	\begin{align*}
		\emph{اتحاديت}&&a + (b+ c) & = (a + b) + c \\
		&& (a \times b) \times c & = a \times (b\times c)\\
		\emph{تبادليت}&&a + b &= b+ a  \\
		&&  a \times b&= b\times a\\
		\emph{عيني غړي}&&a + 0 &= a \\
		&& a \times 1 &= a\\
		\emph{جمعي معکوس}&&(\exists b\in S)0&=a + b\\
		\emph{توزيعيت}&&a \times (b+ c ) &= (a \times b) + (a \times c)
	\end{align*}
	په ضمني ډول دا ټول په هغو کلي مقداري نښو تړل شوي چې ساحه ئې تر $S$ محدوده ده۔ او پام وکړئ چې دلته د $0$ او~$1$ غړي لازماً په همدې نومونو طبيعي عددونه نۀ دي۔
\end{defn}

نو د دې د ارزولو دپاره چې صحيح عددونه تبادلي حلقه جوړوي، يوازې
دا اته شرطونه ارزول غواړو۔ د شرطونو ثابتول {ګران} نۀ دي، خو لږ
اوږد کار دے۔ د بېلګې په ډول د جمع په حالت کښې \emph{اتحاديت}
داسې ثابتوو:

\begin{proof}
$i, j, k \in \Int$ ثابت ونيسئ۔ نو داسې $a_1, b_1, a_2, b_2, a_3, b_3 \in
\Nat$ شته چې $i = \equivrep{a_1, b_1}{}$ او $j =
\equivrep{a_2,b_2}{}$ او $k = \equivrep{a_3, b_3}{}$۔ (د لوستلو
د اسانتيا دپاره د ``$\equivrep{x, y}{\Intequiv}$'' پر ځاے
``$\equivrep{x, y}{}$'' ليکو؛ په دې ټوله برخه کښې به همداسې کوو۔)
اوس:
\begin{align*}
	i + (j + k) &= \equivrep{a_1, b_1}{}+(\equivrep{a_2, b_2}{} + \equivrep{a_3, b_3}{}) \\
	&= \equivrep{a_1,  b_1}{} + \equivrep{a_2+a_3, b_2+b_3}{}\\
	&= \equivrep{a_1 + (a_2 + a_3), b_1 + (b_2 + b_3)}{}\\
	&= \equivrep{(a_1 + a_2) + a_3, (b_1 + b_2) + b_3}{}\\
	&= \equivrep{a_1 + a_2, b_1 + b_2}{} + \equivrep{a_3, b_3}{}\\
	&= (\equivrep{a_1, b_1}{} + \equivrep{a_2, b_2}{}) + \equivrep{a_3, b_3}{}\\
	&= (i+j) + k
\end{align*}
او د $\Nat$ پر جمع له چلنده ازاد کار اخلو۔
\end{proof}

په همدې ډول \emph{جمعي معکوس} داسې ثابتوو:

\begin{proof}
$i \in \Int$ ثابت ونيسئ، نو د ځينو $a,b \in \Nat$ دپاره
$i = \equivrep{a,b}{}$۔ وټاکئ $j = \equivrep{b,a}{} \in \Int$۔
د طبيعي عددونو له چلنده په کار اخيستلو
$(a+b) + 0 = 0 + (a+b)$، نو د تعريف له مخې
$\tuple{a+b, b+a} \sim_\Int \tuple{0,0}$، او له دې امله
$\equivrep{a+b, b+a}{} = \equivrep{0, 0}{} = 0_\Int$۔ نو اوس
$i + j = \equivrep{a,b}{}+\equivrep{b,a}{}=\equivrep{a+b, b+a}{}= \equivrep{0,
0}{} = 0_\Int$۔
\end{proof}

او د \emph{توزيعيت} ثبوت دا دے:

\begin{proof}
لکه پورته، $i = \equivrep{a_1, b_1}{}$،
$j = \equivrep{a_2,b_2}{}$ او $k = \equivrep{a_3, b_3}{}$ ثابت
ونيسئ۔ اوس:
\begin{align*}
	i \times (j + k)
	&= \equivrep{a_1, b_1}{} \times (\equivrep{a_2,b_2}{} + \equivrep{a_3, b_3}{})\\
	&= \equivrep{a_1, b_1}{} \times \equivrep{a_2 + a_3,b_2+b_3}{}\\
	&= \equivrep{a_1  (a_2 + a_3) + b_1  (b_2+b_3), a_1  (b_2 + b_3) + b_1 (a_2 + a_3)}{}\\
	&= \equivrep{a_1 a_2 + a_1a_3 + b_1 b_2+b_1b_3, a_1 b_2 + a_1b_3 + a_2b_1 + a_3b_1}{}\\		
%		&= \equivrep{a_1 a_2 + b_1b_2 + a_1a_3 + b_1 b_2+b_1b_3, a_1 b_2 + a_1b_3 + a_2b_1 + a_3b_1}{}\\		
	&= \equivrep{a_1a_2 + b_1b_2, a_1b_2 + a_2b_1}{} + \equivrep{a_1a_3 + b_1b_3, a_1b_3 + a_3b_1}{}\\
	&= (\equivrep{a_1, b_1}{} \times \equivrep{a_2,b_2}{}) + (\equivrep{a_1, b_1}{} \times  \equivrep{a_3, b_3}{})\\
	&= (i \times j) + (i \times k)
\end{align*}
\end{proof}

د پاتې پنځو شرطونو ثبوت د تمرين په توګه پرېږدو۔ چې دا مو وکړل،
ښودلي مو دي چې $\Int$ يوه تبادلي حلقه جوړوي؛ يعنې جمع او ضرب
(لکه څنګه چې تعريف شوي) هماغسې چلېږي چې بايد۔

\begin{prob}
ثابت کړئ چې $\Int$ يوه تبادلي حلقه ده۔
\end{prob}

خو زموږ کار لا ختم نۀ دے۔ پر $\Int$ مو له جمع او ضرب سره يوه
ترتيبي اړيکه $\leq$ هم تعريف کړې، او بايد وارزوو چې دا لکه څنګه
چې بايد، هماغسې چلېږي۔ په لا تفصيل بايد وښيو چې $\Int$ يوه
\emph{مرتبه} حلقه جوړوي۔\footnote{له
	\olref[sfr][rel][ord]{def:linearorder} څخه په ياد ولرئ چې کلي
	ترتيب هغه اړيکه ده چې انعکاسي، انتقالي، ضدتناظري او وصل شوې وي۔
	د ترتيبي اړيکو په متن کښې وصل والے کله کله \emph{درېګونې} بلل
	کېږي، ځکه د هر $a$ او $b$ دپاره $a \leq b \lor a
	= b \lor a \geq b$ لرو۔}

\begin{defn}
يوه \emph{مرتبه حلقه} هغه تبادلي حلقه ده چې يوه کلي ترتيبي اړيکه
$\leq$ هم لري، داسې چې:
\begin{align*}
	a \leq b &\lif a + c \leq b + c\\
	(a \leq b \land 0 \leq c) &\lif a \times c \leq b \times c
\end{align*}
\end{defn}

\begin{prob}
ثابت کړئ چې $\Int$ يوه مرتبه حلقه ده۔
\end{prob}

لکه مخکښې، دا ښودل چې جوړ شوے~$\Int$ يوه مرتبه حلقه ده، اوږد
خو معمولي کار دے۔ دا تاسو ته پرېږدو۔

په دې سره د صحيح عددونو کار بشپړېږي۔ خو اوس د ناطقو عددونو په اړه
ورته خبرې ښودل غواړو۔ په ځانګړي ډول بايد وښيو چې زموږ د $+$،
$\times$ او $\leq$ د ورکړل شوو تعريفونو له مخې ناطق عددونه يو
مرتب \emph{ميدان} جوړوي:
\begin{defn}\ollabel{orderedfield}
يو \emph{مرتب ميدان} هغه مرتبه حلقه ده چې دا شرط هم پوره کوي:
\begin{align*}
	\emph{ضربي معکوس}& & (\forall a \in S \setminus \{0\})(\exists b \in S) a\times b& = 1
\end{align*}
\textbf{د ژباړې د نسخې بشپړتيا يادونه (OLARI-008):} د ميدان په
دوديز تعريف کښې جمعي او ضربي عيني غړي بايد بېل وي؛ سرچينې دا شرط
نۀ دے بيان کړے، نو ليکل شوے تعريف صفر حلقه هم په تشه توګه مني۔
\end{defn}

چې مو وښودل $\Int$ يوه مرتبه حلقه جوړوي، بيا دا ښودل اسانه خو
اوږد دي چې $\Rat$ يو مرتب ميدان جوړوي۔

\begin{prob}
ثابت کړئ چې $\Rat$ يو مرتب ميدان دے۔
\end{prob}

د صحيح او ناطقو عددونو له سمبالولو وروسته يوازې حقيقي عددونه پاتې
دي۔ په ځانګړي ډول بايد وښيو چې $\Real$ يو \emph{بشپړ} مرتب ميدان
جوړوي، يعنې يو مرتب ميدان چې د بشپړتيا خاصيت لري۔
\olref[cuts]{realcompleteness} ثابته کړه چې $\Real$ د بشپړتيا خاصيت
لري۔ خو لا د دې ارزولو (ستړے) کار پاتې دے چې $\Real$ يو مرتب ميدان
دے۔ \textbf{د ژباړې د نسخې سمون (OLARI-007):} په انګرېزي جمله
کښې د «ستړي» او «ارزولو» تر منځ اړين نوم پاتې شوے و؛ دلته جمله د
همدې کار په څرګند نوم بشپړه شوې ده۔

مخکښې له دې چې \emph{هغه} اوږد تمرين پيل کړو، بايد ځينې نورې
«سمدستي» خبرې وارزوو۔ د بېلګې په ډول بايد ډاډ ولرو چې د هر دوو
پرېکونونو $\alpha$ او $\beta$ دپاره، تعريف شوے $\alpha + \beta$
په رښتيا يو \emph{پرېکون} دے۔ د دې حقيقت ثبوت دا دے:

\begin{proof}	
ځکه چې $\alpha$ او $\beta$ دواړه پرېکونونه دي،
$\alpha + \beta = \Setabs{p + q}{p \in \alpha \land q \in \beta}$
د $\Rat$ يو ناتش خاص فرعي سټ دے۔ اوس فرض کړئ د ځينو
$p \in \alpha$ او $q \in \beta$ دپاره $x < p + q$۔ بيا
$x - p < q$، نو $x - p \in \beta$، او
$x = p + (x - p) \in \alpha + \beta$۔ نو $\alpha + \beta$ د
$\Rat$ يوه پيلنۍ برخه ده۔ په پاے کښې، د هر
$p + q \in \alpha + \beta$ دپاره، ځکه چې $\alpha$ او $\beta$
دواړه پرېکونونه دي، داسې $p_1 \in \alpha$ او $q_1 \in \beta$ شته
چې $p < p_1$ او $q < q_1$؛ نو
$p + q < p_1 + q_1 \in \alpha + \beta$؛ نو
$\alpha + \beta$ تر ټولو لوے غړے نۀ لري۔
\end{proof}

ورته هڅې به درته اجازه وکړي چې وارزوئ $\alpha - \beta$،
$\alpha \times \beta$ او $\alpha \div \beta$ پرېکونونه دي (په
وروستي حالت کښې هغه حالت نۀ شمېرو چې $\beta$ صفر پرېکون وي)۔ خو
يو ځل بيا دا يوازې تاسو ته پرېږدو۔ \textbf{د ژباړې د نسخې بشپړتيا
يادونه (OLARI-009):} د پرېکونونو په برخه کښې د تفريق او وېش
فارمولونه غېر فعالې تبصرې دي، نو فعاله سرچينه هغه دوه عمليې نۀ
تعريفوي چې دا دعوه او راتلونکے ميدان تمرين ورته اړتيا لري۔

\begin{prob}
ثابت کړئ چې $\Real$ يو مرتب ميدان دے۔
\end{prob}

خو يوه وړه نيمګړې خبره لا سمول غواړي۔ په \olref[cuts]{sec} کښې
مو وړانديز وکړ چې $\sqrt{2} =
\Setabs{p \in \Rat}{p < 0 \text{ يا }p^2 < 2}$ اخيستلے شو۔ خو
بايد وښيو چې دا سټ يو \emph{پرېکون} دے۔ د دې حقيقت ثبوت دا دے:

\begin{proof}
په څرګند ډول دا د ناطقو عددونو يوه ناتشه خاصه پيلنۍ برخه ده؛ نو
دومره ښودل بس دي چې تر ټولو لوے غړے نۀ لري۔ په ځانګړي ډول دا ښودل
بس دي چې که $p$ يو مثبت ناطق عدد وي، $p^2 < 2$ وي او
$q = \frac{2p+2}{p+2}$ وي، نو $p < q$ او $q^2 < 2$ دواړه دي۔
د $p < q$ د ليدلو دپاره يوازې پام وکړئ:
\begin{align*}
	p^2 &< 2\\
	p^2 + 2p &< 2 + 2p\\
	p(p + 2) &< 2 + 2p\\
	p &< \tfrac{2+2p}{p+2} = q
\end{align*}
د $q^2 < 2$ د ليدلو دپاره يوازې پام وکړئ:
\begin{align*}
	p^2 &< 2\\
	2p^2 + 4p + 2 &< p^2 + 4p+ 4\\
	4p^2 + 8p + 4 &< 2(p^2 + 4p + 4)\\
	(2p+2)^2 & <2(p+2)^2\\
	\tfrac{(2p+2)^2}{(p+2)^2} &< 2\\
	q^2 &< 2
\end{align*}
\end{proof}
\end{document}
